
\section*{11. Affine group schemes}
\phantomsection\label{sec:sga3-VIB-11}%
\addcontentsline{toc}{section}{11. Affine group schemes}

\medskip
\phantomsection\label{par:sga3-VIB-11.0}%
\noindent\textbf{11.0. Reminders.---}%
\footnote{\label{note:sga3-VIB-104}Editor's note: this
paragraph of reminders has been added.}
Let \(q:X\to S\) be a quasi-compact and quasi-separated
morphism of schemes \textup{(cf.\ EGA~IV\(_1\), 1.1 and
1.2)}, and let \(\mathcal F\) be a quasi-coherent
\(\mathcal O_X\)-module. Recall that \(q_*(\mathcal F)\) is a
quasi-coherent \(\mathcal O_S\)-module
\textup{(EGA~I, 9.2.1)}. Moreover, by EGA~III, 1.4.15
\textup{(completed by EGA~IV\(_1\), 1.7.21)}, one has
\hyperref[concl:sga3-VIB-11.0-c]{\textup{(c)}} below, and the
proof in \emph{loc.\ cit.} also gives
\hyperref[concl:sga3-VIB-11.0-a]{\textup{(a)}} and
\hyperref[concl:sga3-VIB-11.0-b]{\textup{(b)}}:
\begin{enumerate}
\item[\textup{(a)}]\phantomsection\label{concl:sga3-VIB-11.0-a}
  If \(\mathcal F\) is a filtered direct limit of
  quasi-coherent submodules \(\mathcal F_\alpha\), then
  \[
    q_*(\mathcal F)
      =\varinjlim_\alpha q_*(\mathcal F_\alpha).
  \]

\item[\textup{\SGAInlineRefTarget{par:sga3-VIB-11.0:item:b}(b)}]\phantomsection\label{concl:sga3-VIB-11.0-b}
  If \(\mathcal E\) is a flat \(\mathcal O_S\)-module, the
  canonical morphism
  \phantomsection\label{form:sga3-VIB-11.0-projection}%
  \[
    \mathcal E\otimes_{\mathcal O_S}q_*(\mathcal O_X)
      \longrightarrow q_*q^*(\mathcal E)
  \]
  is an isomorphism.

\item[\textup{(c)}]\phantomsection\label{concl:sga3-VIB-11.0-c}
  Let \(p:S'\to S\) be a flat morphism, let
  \(q':X'\to S'\) be obtained from \(q\) by base change, and
  let \(\mathcal F'\) be the inverse image of \(\mathcal F\)
  on \(X'\). Then the canonical morphism
  \phantomsection\label{form:sga3-VIB-11.0-flat-base-change}%
  \[
    p^*q_*(\mathcal F)\longrightarrow q'_*(\mathcal F')
  \]
  is an isomorphism.
\end{enumerate}

Indeed, let \(U=\operatorname{Spec}(A)\) be an arbitrary
affine open subset of \(S\). By hypothesis, \(q^{-1}(U)\) is
the union of affine open subsets
\[
  V_i=\operatorname{Spec}(B_i),\qquad i=1,\ldots,n,
\]
and every intersection \(V_i\cap V_j\) is the union of
finitely many affine open subsets
\[
  W_{ijk}=\operatorname{Spec}(C_{ijk}).
\]
Then
\(\Gamma(U,q_*(\mathcal F))
  =\Gamma(q^{-1}(U),\mathcal F)\)
is the kernel of the morphism
\phantomsection\label{form:sga3-VIB-11.0-Cech-F}%
\[
  \bigoplus_{i=1}^{n}\Gamma(V_i,\mathcal F)
    \longrightarrow
  \bigoplus_{i,j,k}\Gamma(W_{ijk},\mathcal F).
\]
Assertion~\hyperref[concl:sga3-VIB-11.0-a]{\textup{(a)}}
follows, since every term above commutes with filtered direct
limits \textup{(the \(V_i\) and \(W_{ijk}\) being affine,
hence quasi-compact)}.

Let us prove
\hyperref[concl:sga3-VIB-11.0-b]{\textup{(b)}}. The module
\(E=\Gamma(U,\mathcal E)\) is flat over \(A\), and
\(\Gamma(U,q_*q^*(\mathcal E))\) is the kernel \(K(E)\) of
the morphism
\phantomsection\label{form:sga3-VIB-11.0-Cech-E}%
\[
  \bigoplus_{i=1}^{n}B_i\otimes_AE
    \longrightarrow
  \bigoplus_{i,j,k}C_{ijk}\otimes_AE.
\]
Since \(E\) is flat over \(A\), this kernel is identified
with
\[
  K(A)\otimes_AE
    =\mathcal O_X(q^{-1}(U))\otimes_AE.
\]
Finally, if \(U'\) is an arbitrary affine open subset of
\(S'\) lying over \(U\), then
\(A'=\mathcal O(U')\) is a flat \(A\)-algebra, and the same
argument gives%
\begingroup
\renewcommand{\thefootnote}{\fnsymbol{footnote}}%
\footnote[1]{\label{note:sga3-VIB-11.0-U-prime}
Translator's note: the French re-edition prints
\(q'^{-1}(U)\) in the right-hand side below. Since
\(q':X'\to S'\) and the affine open subset just introduced
is \(U'\subset S'\), the inverse image is
\(q'^{-1}(U')\).}%
\endgroup
\phantomsection\label{form:sga3-VIB-11.0-local-base-change}%
\[
  \mathcal F(q^{-1}(U))\otimes_AA'
    \xrightarrow{\ \sim\ }
  \mathcal F'(q'^{-1}(U')).
\]

\medskip
\phantomsection\label{not:sga3-VIB-11-A}%
\noindent\textbf{Notation.---}
Let \(S\) be a scheme, let \(X\) be an \(S\)-scheme, and let
\(f:X\to S\) be the structural morphism. We shall put
\phantomsection\label{form:sga3-VIB-11-A-X}%
\[
  \mathcal A(X)=f_*(\mathcal O_X).
\]

\medskip
\phantomsection\label{lem:sga3-VIB-11.1}%
\noindent\textbf{Lemma 11.1.---}
\emph{Let \(X\) and \(Y\) be two quasi-compact and
quasi-separated \(S\)-schemes, and let \(f:X\to S\) and
\(g:Y\to S\) be their structural morphisms. Then the
canonical homomorphism}
\phantomsection\label{form:sga3-VIB-11.1-phi}%
\[
  \varphi:
  \mathcal A(X)\otimes_{\mathcal O_S}\mathcal A(Y)
    \longrightarrow
  \mathcal A(X\times_SY)
\]
\emph{is an isomorphism in each of the following cases:}%
\footnote{\label{note:sga3-VIB-105}Editor's note: if \(k\) is
a field and \(X\) is an infinite disjoint union of copies of
\(S=\operatorname{Spec}k\) \textup{(so that \(X\) is not
quasi-compact)}, then \(\mathcal A(X)=k^X\), and the
canonical morphism
\[
  k^X\otimes_k k^X\longrightarrow k^{X\times X}
\]
is not surjective.}
\begin{enumerate}
\item[\textup{(a)}]\phantomsection\label{concl:sga3-VIB-11.1-a}
  \emph{\(f\) and \(g\) are affine;}

\item[\textup{(b)}]\phantomsection\label{concl:sga3-VIB-11.1-b}
  \emph{\(g\) \textup{(or \(f\))} is flat and affine;}

\item[\textup{(c)}]\phantomsection\label{concl:sga3-VIB-11.1-c}
  \emph{\(g\) is flat and \(f_*(\mathcal O_X)\) is a flat
  \(\mathcal O_S\)-module.}
\end{enumerate}

In case~\hyperref[concl:sga3-VIB-11.1-b]{\textup{(b)}}, we
shall suppose that it is \(g\) that is flat and affine. Put
\[
  S'=\operatorname{Spec}\mathcal A(X),\qquad
  Y'=Y\times_SS',\qquad
  g'=g\times_SS',
\]
and write \(v:S'\to S\) for the structural morphism:
\begin{center}
  \phantomsection\label{diag:sga3-VIB-11.1-base-change}%
  \begin{tikzcd}[row sep=large,column sep=large]
    Y' \arrow[r] \arrow[d,"g'"'] &
    Y \arrow[d,"g"] \\
    \operatorname{Spec}\mathcal A(X)=S'
      \arrow[r,"v"'] &
    S
  \end{tikzcd}
\end{center}
In cases~\hyperref[concl:sga3-VIB-11.1-a]{\textup{(a)}} and
\hyperref[concl:sga3-VIB-11.1-b]{\textup{(b)}}, the morphism
\(g\) is affine, and hence EGA~II, 1.5.2 gives
\phantomsection\label{form:sga3-VIB-11.1-eq1}%
\[
  \tag{1}
  g'_*(\mathcal O_{Y'})
    =v^*g_*(\mathcal O_Y)
    =\mathcal A(Y)\otimes_{\mathcal O_S}\mathcal O_{S'}.
\]
The same equality holds in
case~\hyperref[concl:sga3-VIB-11.1-c]{\textup{(c)}}, by
\hyperref[concl:sga3-VIB-11.0-c]{11.0\textup{(c)}}, since
\(S'\) is flat over \(S\) and \(g\) is quasi-compact and
quasi-separated.

On the other hand, by EGA~II, 1.2.7, the morphism
\(f:X\to S\) factors through \(v\) by means of a morphism
\(p:X\to S'\), and
\[
  p_*(\mathcal O_X)=\mathcal O_{S'},
  \qquad
  X\times_SY=X\times_{S'}Y'.
\]
Since \(f\) is quasi-separated, so is \(p\)
\textup{(EGA~IV\(_1\), 1.2.2)}; and since \(f\) is
quasi-compact and \(v\) is quasi-separated, \(p\) is also
quasi-compact \textup{(EGA~IV\(_1\), 1.2.4)}. Consider the
cartesian square
\begin{center}
  \phantomsection\label{diag:sga3-VIB-11.1-cartesian}%
  \begin{tikzcd}[row sep=large,column sep=large]
    X\times_SY \arrow[r,"p'"] \arrow[d] &
    Y' \arrow[d,"g'"] \\
    X \arrow[r,"p"'] &
    S'
  \end{tikzcd}
\end{center}
In cases~\hyperref[concl:sga3-VIB-11.1-b]{\textup{(b)}} and
\hyperref[concl:sga3-VIB-11.1-c]{\textup{(c)}}, the scheme
\(Y\) is flat over \(S\), and hence \(Y'\) is flat over
\(S'\). Applying
\hyperref[concl:sga3-VIB-11.0-c]{11.0\textup{(c)}} once
again, we obtain
\phantomsection\label{form:sga3-VIB-11.1-eq2}%
\[
  \tag{2}
  p'_*(\mathcal O_{X\times_SY})
    =g'^*p_*(\mathcal O_X)
    =g'^*(\mathcal O_{S'})
    =\mathcal O_{Y'}.
\]
The same equality holds in
case~\hyperref[concl:sga3-VIB-11.1-a]{\textup{(a)}}, because
\(p\) and \(p'\) are then isomorphisms.

Finally, since \(v\) is affine, EGA~II, 1.4.7 gives
\[
  v_*(\mathcal F\otimes_{\mathcal O_S}\mathcal O_{S'})
    =\mathcal F\otimes_{\mathcal O_S}\mathcal A(X)
\]
for every quasi-coherent \(\mathcal O_S\)-module
\(\mathcal F\). Combined with
\hyperref[form:sga3-VIB-11.1-eq2]{\textup{(2)}} and
\hyperref[form:sga3-VIB-11.1-eq1]{\textup{(1)}}, this gives
\[
\begin{aligned}
  \mathcal A(X\times_SY)
    &=v_*g'_*p'_*(\mathcal O_{X\times_SY})\\
    &\overset{\textup{(2)}}{=}v_*g'_*(\mathcal O_{Y'})
     \overset{\textup{(1)}}{=}
       v_*(\mathcal A(Y)\otimes_{\mathcal O_S}\mathcal O_{S'})\\
    &=\mathcal A(Y)\otimes_{\mathcal O_S}\mathcal A(X).
\end{aligned}
\]

\medskip
\phantomsection\label{cor:sga3-VIB-11.2}%
\noindent\textbf{Corollary 11.2.---}
\emph{The functor}
\[
  X\longmapsto\operatorname{Spec}\mathcal A(X),
\]
\emph{from the full subcategory of
\((\mathbf{Sch}/S)\) formed by the \(S\)-schemes \(X\) that
are flat, quasi-compact, and quasi-separated over \(S\), and
for which \(\mathcal A(X)\) is a flat
\(\mathcal O_S\)-module, to the category of flat affine
\(S\)-schemes, commutes with finite products and therefore
takes \(S\)-groups to \(S\)-groups.}

\medskip
\phantomsection\label{def:sga3-VIB-11.3}%
\noindent\textbf{Definition 11.3.---}
\emph{Given a flat \(S\)-group \(G\), quasi-compact and
quasi-separated over \(S\), such that \(\mathcal A(G)\) is
flat over \(\mathcal O_S\),}%
\footnote{\label{note:sga3-VIB-106}Editor's note: if \(S\) is
a regular locally Noetherian scheme of dimension at most
\(2\), and \(X\) is an \(S\)-scheme that is flat,
quasi-compact, and quasi-separated, then \(\mathcal A(X)\) is
a flat \(\mathcal O_S\)-module; cf.\ \textup{[Ray70a]},
VII, \SGARefLink{sga3:VIIB:3.2:prop}{3.2}.}
\emph{we write \(G_{\mathrm{af}}\), and call the
\(S\)-group}
\[
  G_{\mathrm{af}}=\operatorname{Spec}\mathcal A(G)
\]
\emph{the \emph{affine envelope} of \(G\).}

\medskip
\phantomsection\label{prop:sga3-VIB-11.3.1}%
\noindent\textbf{Proposition 11.3.1.---}%
\footnote{\label{note:sga3-VIB-107}Editor's note: this
proposition has been added; see also the additional
\hyperref[sec:sga3-VIB-12]{Section~12} below for a study of the morphism
\(G\to G_{\mathrm{af}}\) and of its kernel.}
\emph{The canonical morphism}
\[
  \tau_G:G\longrightarrow G_{\mathrm{af}}
\]
\emph{is a morphism of \(S\)-groups. Moreover, it has the
following universal property:}
\begin{enumerate}
\item[\textup{(i)}]\phantomsection\label{concl:sga3-VIB-11.3.1-i}
  \emph{For every morphism of \(S\)-schemes
  \(\varphi:G\to H\), where \(H\) is affine over \(S\), there
  is a unique morphism of \(S\)-schemes
  \(\varphi':G_{\mathrm{af}}\to H\) such that}
  \[
    \varphi=\varphi'\circ\tau_G.
  \]

\item[\textup{(ii)}]\phantomsection\label{concl:sga3-VIB-11.3.1-ii}
  \emph{If, moreover, \(H\) is an \(S\)-group and
  \(\varphi\) is a morphism of \(S\)-groups, then the same is
  true of \(\varphi'\).}
\end{enumerate}
